% Options for packages loaded elsewhere
\PassOptionsToPackage{unicode}{hyperref}
\PassOptionsToPackage{hyphens}{url}
\PassOptionsToPackage{dvipsnames,svgnames,x11names}{xcolor}
%
\documentclass[
  10pt,
]{article}
\usepackage{amsmath,amssymb}
\usepackage{iftex}
\ifPDFTeX
  \usepackage[T1]{fontenc}
  \usepackage[utf8]{inputenc}
  \usepackage{textcomp} % provide euro and other symbols
\else % if luatex or xetex
  \usepackage{unicode-math} % this also loads fontspec
  \defaultfontfeatures{Scale=MatchLowercase}
  \defaultfontfeatures[\rmfamily]{Ligatures=TeX,Scale=1}
\fi
\usepackage{lmodern}
\ifPDFTeX\else
  % xetex/luatex font selection
\fi
% Use upquote if available, for straight quotes in verbatim environments
\IfFileExists{upquote.sty}{\usepackage{upquote}}{}
\IfFileExists{microtype.sty}{% use microtype if available
  \usepackage[]{microtype}
  \UseMicrotypeSet[protrusion]{basicmath} % disable protrusion for tt fonts
}{}
\makeatletter
\@ifundefined{KOMAClassName}{% if non-KOMA class
  \IfFileExists{parskip.sty}{%
    \usepackage{parskip}
  }{% else
    \setlength{\parindent}{0pt}
    \setlength{\parskip}{6pt plus 2pt minus 1pt}}
}{% if KOMA class
  \KOMAoptions{parskip=half}}
\makeatother
\usepackage{xcolor}
\usepackage[margin=2.2cm]{geometry}
\usepackage{color}
\usepackage{fancyvrb}
\newcommand{\VerbBar}{|}
\newcommand{\VERB}{\Verb[commandchars=\\\{\}]}
\DefineVerbatimEnvironment{Highlighting}{Verbatim}{commandchars=\\\{\}}
% Add ',fontsize=\small' for more characters per line
\newenvironment{Shaded}{}{}
\newcommand{\AlertTok}[1]{\textcolor[rgb]{1.00,0.00,0.00}{\textbf{#1}}}
\newcommand{\AnnotationTok}[1]{\textcolor[rgb]{0.38,0.63,0.69}{\textbf{\textit{#1}}}}
\newcommand{\AttributeTok}[1]{\textcolor[rgb]{0.49,0.56,0.16}{#1}}
\newcommand{\BaseNTok}[1]{\textcolor[rgb]{0.25,0.63,0.44}{#1}}
\newcommand{\BuiltInTok}[1]{\textcolor[rgb]{0.00,0.50,0.00}{#1}}
\newcommand{\CharTok}[1]{\textcolor[rgb]{0.25,0.44,0.63}{#1}}
\newcommand{\CommentTok}[1]{\textcolor[rgb]{0.38,0.63,0.69}{\textit{#1}}}
\newcommand{\CommentVarTok}[1]{\textcolor[rgb]{0.38,0.63,0.69}{\textbf{\textit{#1}}}}
\newcommand{\ConstantTok}[1]{\textcolor[rgb]{0.53,0.00,0.00}{#1}}
\newcommand{\ControlFlowTok}[1]{\textcolor[rgb]{0.00,0.44,0.13}{\textbf{#1}}}
\newcommand{\DataTypeTok}[1]{\textcolor[rgb]{0.56,0.13,0.00}{#1}}
\newcommand{\DecValTok}[1]{\textcolor[rgb]{0.25,0.63,0.44}{#1}}
\newcommand{\DocumentationTok}[1]{\textcolor[rgb]{0.73,0.13,0.13}{\textit{#1}}}
\newcommand{\ErrorTok}[1]{\textcolor[rgb]{1.00,0.00,0.00}{\textbf{#1}}}
\newcommand{\ExtensionTok}[1]{#1}
\newcommand{\FloatTok}[1]{\textcolor[rgb]{0.25,0.63,0.44}{#1}}
\newcommand{\FunctionTok}[1]{\textcolor[rgb]{0.02,0.16,0.49}{#1}}
\newcommand{\ImportTok}[1]{\textcolor[rgb]{0.00,0.50,0.00}{\textbf{#1}}}
\newcommand{\InformationTok}[1]{\textcolor[rgb]{0.38,0.63,0.69}{\textbf{\textit{#1}}}}
\newcommand{\KeywordTok}[1]{\textcolor[rgb]{0.00,0.44,0.13}{\textbf{#1}}}
\newcommand{\NormalTok}[1]{#1}
\newcommand{\OperatorTok}[1]{\textcolor[rgb]{0.40,0.40,0.40}{#1}}
\newcommand{\OtherTok}[1]{\textcolor[rgb]{0.00,0.44,0.13}{#1}}
\newcommand{\PreprocessorTok}[1]{\textcolor[rgb]{0.74,0.48,0.00}{#1}}
\newcommand{\RegionMarkerTok}[1]{#1}
\newcommand{\SpecialCharTok}[1]{\textcolor[rgb]{0.25,0.44,0.63}{#1}}
\newcommand{\SpecialStringTok}[1]{\textcolor[rgb]{0.73,0.40,0.53}{#1}}
\newcommand{\StringTok}[1]{\textcolor[rgb]{0.25,0.44,0.63}{#1}}
\newcommand{\VariableTok}[1]{\textcolor[rgb]{0.10,0.09,0.49}{#1}}
\newcommand{\VerbatimStringTok}[1]{\textcolor[rgb]{0.25,0.44,0.63}{#1}}
\newcommand{\WarningTok}[1]{\textcolor[rgb]{0.38,0.63,0.69}{\textbf{\textit{#1}}}}
\usepackage{longtable,booktabs,array}
\usepackage{calc} % for calculating minipage widths
% Correct order of tables after \paragraph or \subparagraph
\usepackage{etoolbox}
\makeatletter
\patchcmd\longtable{\par}{\if@noskipsec\mbox{}\fi\par}{}{}
\makeatother
% Allow footnotes in longtable head/foot
\IfFileExists{footnotehyper.sty}{\usepackage{footnotehyper}}{\usepackage{footnote}}
\makesavenoteenv{longtable}
\setlength{\emergencystretch}{3em} % prevent overfull lines
\providecommand{\tightlist}{%
  \setlength{\itemsep}{0pt}\setlength{\parskip}{0pt}}
\setcounter{secnumdepth}{-\maxdimen} % remove section numbering
\ifLuaTeX
\usepackage[bidi=basic]{babel}
\else
\usepackage[bidi=default]{babel}
\fi
\babelprovide[main,import]{italian}
% get rid of language-specific shorthands (see #6817):
\let\LanguageShortHands\languageshorthands
\def\languageshorthands#1{}
\usepackage{needspace}
\usepackage{etoolbox}
\pretocmd{\subsection}{\Needspace{12\baselineskip}}{}{}
\pretocmd{\section}{\Needspace{14\baselineskip}}{}{}
\ifLuaTeX
  \usepackage{selnolig}  % disable illegal ligatures
\fi
\usepackage{bookmark}
\IfFileExists{xurl.sty}{\usepackage{xurl}}{} % add URL line breaks if available
\urlstyle{same}
\hypersetup{
  pdflang={it-IT},
  colorlinks=true,
  linkcolor={Maroon},
  filecolor={Maroon},
  citecolor={Blue},
  urlcolor={Blue},
  pdfcreator={LaTeX via pandoc}}

\author{}
\date{}

\begin{document}

\section{Facial-access Decision demo - independent authoring
corpus}\label{facial-access-decision-demo---independent-authoring-corpus}

\begin{quote}
\textbf{Non-canonical research example - Phase D1.}

This corpus was authored after commit
\texttt{4fd39955a1920e4380c0dc411b0d814c433646a6}, using the frozen
Decision construction candidate and the four previously studied
facial-access Macro Requirements. It is intentionally written
\textbf{before} Phase D2 analysis. No statement in this document is
evidence that the candidate metamodel fits; that judgement is deferred
to a separate checkpoint.
\end{quote}

\subsection{Authoring boundary}\label{authoring-boundary}

The Decision-specific shape is kept fixed during this phase:

\begin{Shaded}
\begin{Highlighting}[]
\NormalTok{Decision}
\NormalTok{|{-} title}
\NormalTok{|{-} context}
\NormalTok{|{-} decision}
\NormalTok{\textasciigrave{}{-} consequences}
\end{Highlighting}
\end{Shaded}

Every Decision below materializes all four fields. Parent Macro
Requirement is recorded as an explicit relation in the study wrapper,
not duplicated as a Decision-specific semantic field. Lifecycle/status,
Requirement semantics, architecture topology, concrete protocols,
thresholds, database choices and deployment placement remain outside
this authoring exercise.

The independent source MRs are:

\begin{enumerate}
\def\labelenumi{\arabic{enumi}.}
\tightlist
\item
  \texttt{MR-0001} - Accesso controllato all'area riservata
\item
  \texttt{MR-0002} - Gestione delle persone autorizzate
\item
  \texttt{MR-0003} - Riconoscimento facciale per la verifica della
  persona
\item
  \texttt{MR-0004} - Uso responsabile dei dati biometrici e delle
  decisioni automatiche
\end{enumerate}

\subsection{Corpus overview}\label{corpus-overview}

\begin{longtable}[]{@{}
  >{\raggedright\arraybackslash}p{(\columnwidth - 4\tabcolsep) * \real{0.3333}}
  >{\raggedright\arraybackslash}p{(\columnwidth - 4\tabcolsep) * \real{0.3333}}
  >{\raggedright\arraybackslash}p{(\columnwidth - 4\tabcolsep) * \real{0.3333}}@{}}
\toprule\noalign{}
\begin{minipage}[b]{\linewidth}\raggedright
Parent MR
\end{minipage} & \begin{minipage}[b]{\linewidth}\raggedright
Decision
\end{minipage} & \begin{minipage}[b]{\linewidth}\raggedright
Title
\end{minipage} \\
\midrule\noalign{}
\endhead
\bottomrule\noalign{}
\endlastfoot
MR-0001 & ADR-0001 & Autorizzazione e verifica come condizioni distinte
dell'apertura automatica \\
MR-0001 & ADR-0002 & Un esito non conclusivo non equivale a verifica
riuscita \\
MR-0002 & ADR-0001 & Autorizzazione organizzativa separata
dall'iscrizione biometrica \\
MR-0002 & ADR-0002 & Una fonte autorevole per lo stato corrente di
autorizzazione \\
MR-0003 & ADR-0001 & Il riconoscimento produce evidenza di identita, non
una decisione di accesso \\
MR-0003 & ADR-0002 & Esito di riconoscimento semanticamente stabile
rispetto al modello ML \\
MR-0004 & ADR-0001 & Limitazione di finalita dei dati biometrici \\
MR-0004 & ADR-0002 & Decisioni automatiche contestabili e
revisionabili \\
\end{longtable}

\section{MR-0001 - Accesso controllato all'area
riservata}\label{mr-0001---accesso-controllato-allarea-riservata}

\subsection{ADR-0001 - Autorizzazione e verifica come condizioni
distinte dell'apertura
automatica}\label{adr-0001---autorizzazione-e-verifica-come-condizioni-distinte-dellapertura-automatica}

\textbf{Parent Macro Requirement:} \texttt{MR-0001}

\subsubsection{Title}\label{title}

Autorizzazione e verifica come condizioni distinte dell'apertura
automatica

\subsubsection{Context}\label{context}

L'apertura automatica del varco deve riflettere sia la decisione
organizzativa su chi e autorizzato sia l'evidenza che la persona
presente al punto di accesso corrisponda all'identita interessata. Se il
risultato del riconoscimento facciale viene trattato direttamente come
autorizzazione, la capacita biometrica finisce per sostituire una
decisione organizzativa che ha significato autonomo e puo cambiare nel
tempo.

\subsubsection{Decision}\label{decision}

Il progetto tratta \textbf{autorizzazione corrente} e \textbf{verifica
della persona} come condizioni distinte della decisione di apertura
automatica. Il riconoscimento facciale puo fornire evidenza per la
verifica della persona, ma non crea, estende o sostituisce
l'autorizzazione. L'apertura automatica rappresenta quindi un esito del
processo di accesso soltanto quando le condizioni di autorizzazione e
verifica richieste dal progetto risultano entrambe soddisfatte.

\subsubsection{Consequences}\label{consequences}

\begin{itemize}
\tightlist
\item
  La responsabilita di stabilire chi puo accedere resta distinta dalla
  tecnologia usata per riconoscere la persona.
\item
  Un riconoscimento positivo non puo, da solo, rendere autorizzata una
  persona.
\item
  Una persona autorizzata puo non ottenere apertura automatica quando la
  verifica richiesta non e disponibile o sufficiente.
\item
  I livelli successivi dovranno rendere osservabile come le due
  condizioni vengono risolte e combinate senza duplicarne le fonti
  autorevoli.
\end{itemize}

\subsection{ADR-0002 - Un esito non conclusivo non equivale a verifica
riuscita}\label{adr-0002---un-esito-non-conclusivo-non-equivale-a-verifica-riuscita}

\textbf{Parent Macro Requirement:} \texttt{MR-0001}

\subsubsection{Title}\label{title-1}

Un esito non conclusivo non equivale a verifica riuscita

\subsubsection{Context}\label{context-1}

Il riconoscimento facciale puo non produrre evidenza sufficiente per
stabilire una corrispondenza affidabile. Trattare un esito incerto,
incompleto o non disponibile come se fosse una verifica positiva
renderebbe l'apertura automatica dipendente dall'assenza di evidenza
anziche dalla sua presenza; trattarlo sempre come un rifiuto definitivo
confonderebbe invece l'impossibilita di verificare con la certezza che
la persona non corrisponda.

\subsubsection{Decision}\label{decision-1}

Il processo di accesso conserva l'\textbf{esito non conclusivo} come
situazione distinta dalla verifica riuscita e dalla verifica negativa.
Un esito non conclusivo non vale come verifica positiva ai fini
dell'apertura automatica. Il modo con cui il progetto gestira
successivamente l'eccezione, la ripetizione della verifica o un
eventuale intervento organizzativo rimane una scelta successiva.

\subsubsection{Consequences}\label{consequences-1}

\begin{itemize}
\tightlist
\item
  L'apertura automatica non viene concessa sulla base di un risultato
  che il sistema non e in grado di interpretare come verifica riuscita.
\item
  Il progetto preserva la differenza tra mancata corrispondenza e
  impossibilita di concludere.
\item
  L'esperienza di accesso puo richiedere percorsi aggiuntivi nei casi
  non conclusivi.
\item
  Le scelte successive dovranno mantenere questa distinzione senza
  trasformarla prematuramente in una specifica tecnologia o procedura.
\end{itemize}

\section{MR-0002 - Gestione delle persone
autorizzate}\label{mr-0002---gestione-delle-persone-autorizzate}

\subsection{ADR-0001 - Autorizzazione organizzativa separata
dall'iscrizione
biometrica}\label{adr-0001---autorizzazione-organizzativa-separata-dalliscrizione-biometrica}

\textbf{Parent Macro Requirement:} \texttt{MR-0002}

\subsubsection{Title}\label{title-2}

Autorizzazione organizzativa separata dall'iscrizione biometrica

\subsubsection{Context}\label{context-2}

Per riconoscere una persona il progetto deve poter associare alla sua
identita informazioni biometriche adeguate. La presenza di tali
informazioni, tuttavia, non significa necessariamente che la persona sia
attualmente autorizzata ad accedere: autorizzazioni e dati necessari al
riconoscimento hanno motivazioni e cicli di cambiamento differenti.

\subsubsection{Decision}\label{decision-2}

Il progetto mantiene \textbf{l'autorizzazione organizzativa} e
\textbf{l'iscrizione biometrica necessaria alla verifica} come fatti
distinti. L'iscrizione biometrica non concede di per se accesso e una
modifica dello stato di autorizzazione non viene rappresentata come
modifica implicita dell'identita della persona. Eventuali regole che
coordinano i rispettivi cicli di vita dovranno essere governate
esplicitamente a un livello successivo.

\subsubsection{Consequences}\label{consequences-2}

\begin{itemize}
\tightlist
\item
  La concessione o revoca dell'accesso non dipende dal semplice fatto
  che esista materiale biometrico utilizzabile.
\item
  Il progetto puo ragionare separatamente sulla correttezza
  dell'autorizzazione e sulla qualita delle informazioni di
  riconoscimento.
\item
  La separazione rende visibili future domande su conservazione,
  aggiornamento e cancellazione dei dati biometrici invece di
  nasconderle nello stato di accesso.
\item
  I livelli successivi dovranno mantenere coerenza tra fatti distinti
  senza trasformarli in copie concorrenti della stessa informazione.
\end{itemize}

\subsection{ADR-0002 - Una fonte autorevole per lo stato corrente di
autorizzazione}\label{adr-0002---una-fonte-autorevole-per-lo-stato-corrente-di-autorizzazione}

\textbf{Parent Macro Requirement:} \texttt{MR-0002}

\subsubsection{Title}\label{title-3}

Una fonte autorevole per lo stato corrente di autorizzazione

\subsubsection{Context}\label{context-3}

Lo stato di autorizzazione puo cambiare nel tempo e viene consumato dal
processo che decide l'accesso. Se lo stesso stato viene mantenuto come
verita indipendente in piu punti del progetto, una sospensione o revoca
puo risultare applicata in modo non uniforme e il sistema puo continuare
a operare su copie obsolete.

\subsubsection{Decision}\label{decision-3}

Il progetto mantiene \textbf{una sola fonte autorevole dello stato
corrente di autorizzazione} per ciascuna identita governata. Le
rappresentazioni necessarie ai diversi consumatori sono derivate o
risolte da quella fonte e non costituiscono nuove autorita indipendenti
sul significato dell'autorizzazione.

\subsubsection{Consequences}\label{consequences-3}

\begin{itemize}
\tightlist
\item
  La decisione organizzativa corrente rimane ricostruibile senza
  confrontare copie concorrenti.
\item
  Proiezioni o copie operative possono esistere, ma la loro correttezza
  dipende dalla relazione con la fonte autorevole.
\item
  Disponibilita, aggiornamento e freshness delle rappresentazioni
  derivate diventano problemi espliciti da governare successivamente.
\item
  La scelta limita soluzioni che attribuiscono autonomia semantica a uno
  stato locale vicino al componente di riconoscimento o al varco.
\end{itemize}

\section{MR-0003 - Riconoscimento facciale per la verifica della
persona}\label{mr-0003---riconoscimento-facciale-per-la-verifica-della-persona}

\subsection{ADR-0001 - Il riconoscimento produce evidenza di identita,
non una decisione di
accesso}\label{adr-0001---il-riconoscimento-produce-evidenza-di-identita-non-una-decisione-di-accesso}

\textbf{Parent Macro Requirement:} \texttt{MR-0003}

\subsubsection{Title}\label{title-4}

Il riconoscimento produce evidenza di identita, non una decisione di
accesso

\subsubsection{Context}\label{context-4}

Un sistema di riconoscimento facciale puo essere progettato per
restituire direttamente etichette operative come ``accesso consentito''
oppure per descrivere cio che ha osservato rispetto a un'identita.
Legare il significato dell'output del riconoscimento alla politica di
accesso renderebbe il modello biometrico proprietario anche di decisioni
che dipendono da autorizzazioni e regole esterne al riconoscimento
stesso.

\subsubsection{Decision}\label{decision-4}

La capacita di riconoscimento facciale produce \textbf{evidenza sulla
corrispondenza tra il volto osservato e un'identita governata}, non una
decisione autonoma di accesso. Il processo di accesso interpreta tale
evidenza insieme alle altre conoscenze governate necessarie alla propria
decisione.

\subsubsection{Consequences}\label{consequences-4}

\begin{itemize}
\tightlist
\item
  Il modello di riconoscimento resta separato dalla politica che
  stabilisce chi puo entrare.
\item
  Cambiare una regola di autorizzazione non richiede di ridefinire il
  significato del riconoscimento facciale.
\item
  Il processo di accesso deve integrare esplicitamente l'evidenza di
  riconoscimento con le altre condizioni rilevanti.
\item
  Il risultato del riconoscimento deve restare sufficientemente
  comprensibile da poter essere valutato e riesaminato senza dipendere
  da una singola etichetta operativa.
\end{itemize}

\subsection{ADR-0002 - Esito di riconoscimento semanticamente stabile
rispetto al modello
ML}\label{adr-0002---esito-di-riconoscimento-semanticamente-stabile-rispetto-al-modello-ml}

\textbf{Parent Macro Requirement:} \texttt{MR-0003}

\subsubsection{Title}\label{title-5}

Esito di riconoscimento semanticamente stabile rispetto al modello ML

\subsubsection{Context}\label{context-5}

Modelli o implementazioni differenti possono esporre punteggi, scale e
indicatori con significati tecnici diversi. Se il resto del progetto
dipende direttamente da un valore grezzo specifico del modello,
sostituire o aggiornare il modello puo cambiare implicitamente il
significato della verifica anche quando l'intento del progetto rimane
invariato.

\subsubsection{Decision}\label{decision-5}

Il progetto definisce un \textbf{significato stabile dell'esito di
riconoscimento} a livello della capacita di verifica, separato dalla
scala o dal formato grezzo prodotti da uno specifico modello ML.
Evidenze tecniche del modello possono essere conservate per spiegazione,
revisione e valutazione, ma il contratto semantico consumato dal
processo di accesso non coincide con un punteggio grezzo proprietario di
una particolare implementazione.

\subsubsection{Consequences}\label{consequences-5}

\begin{itemize}
\tightlist
\item
  L'evoluzione del modello ML puo essere valutata senza cambiare
  automaticamente la semantica del processo di accesso.
\item
  Diventa necessario governare come l'evidenza tecnica viene
  interpretata nel significato stabile della verifica.
\item
  Informazioni specifiche del modello restano disponibili come evidenza
  senza diventare l'unico linguaggio condiviso dal progetto.
\item
  La scelta introduce un confine esplicito tra semantica del progetto e
  dettagli dell'implementazione ML.
\end{itemize}

\section{MR-0004 - Uso responsabile dei dati biometrici e delle
decisioni
automatiche}\label{mr-0004---uso-responsabile-dei-dati-biometrici-e-delle-decisioni-automatiche}

\subsection{ADR-0001 - Limitazione di finalita dei dati
biometrici}\label{adr-0001---limitazione-di-finalita-dei-dati-biometrici}

\textbf{Parent Macro Requirement:} \texttt{MR-0004}

\subsubsection{Title}\label{title-6}

Limitazione di finalita dei dati biometrici

\subsubsection{Context}\label{context-6}

Le informazioni biometriche raccolte per consentire il riconoscimento al
varco possono tecnicamente essere riutilizzate per finalita ulteriori,
come osservazione continuativa, profilazione o analisi non necessarie al
controllo dell'accesso. Un riuso implicito allargherebbe il significato
del progetto senza una corrispondente scelta governata e senza che le
persone coinvolte possano comprenderne il nuovo scopo.

\subsubsection{Decision}\label{decision-6}

Il progetto usa i dati biometrici per \textbf{iscrizione, verifica della
persona e revisione controllata direttamente connesse alla finalita di
accesso}. I dati non vengono trattati come una risorsa genericamente
disponibile per finalita estranee al controllo dell'accesso. Una
finalita sostanzialmente diversa richiederebbe una nuova scelta
governata e non viene assorbita implicitamente da questa Decision.

\subsubsection{Consequences}\label{consequences-6}

\begin{itemize}
\tightlist
\item
  Le capacita future non possono presumere che i dati biometrici siano
  disponibili per qualunque uso tecnicamente possibile.
\item
  La comprensibilita dello scopo del trattamento migliora per persone e
  responsabili del progetto.
\item
  Alcune funzionalita secondarie possono richiedere una nuova
  giustificazione e una nuova decomposizione del progetto invece di
  essere aggiunte come semplice riuso dei dati.
\item
  I livelli successivi dovranno rendere operativa la limitazione di
  finalita senza duplicarne manualmente il significato in ogni documento
  o componente.
\end{itemize}

\subsection{ADR-0002 - Decisioni automatiche contestabili e
revisionabili}\label{adr-0002---decisioni-automatiche-contestabili-e-revisionabili}

\textbf{Parent Macro Requirement:} \texttt{MR-0004}

\subsubsection{Title}\label{title-7}

Decisioni automatiche contestabili e revisionabili

\subsubsection{Context}\label{context-7}

Il riconoscimento facciale puo produrre errori o risultati non
conclusivi che incidono sull'accesso fisico di una persona. Se l'esito
automatico fosse trattato come autorita finale non riesaminabile, un
errore tecnico diventerebbe direttamente una decisione organizzativa
senza un percorso con cui comprenderne o contestarne gli effetti.

\subsubsection{Decision}\label{decision-7}

Gli esiti automatici che incidono negativamente sull'accesso restano
\textbf{contestabili e revisionabili attraverso un processo
organizzativo distinto dalla sola esecuzione del modello}. L'output
automatico costituisce evidenza per la gestione dell'accesso, ma non
elimina la possibilita di una valutazione successiva da parte
dell'organizzazione.

\subsubsection{Consequences}\label{consequences-7}

\begin{itemize}
\tightlist
\item
  Le persone dispongono di un principio di riesame quando l'automazione
  produce un esito problematico.
\item
  L'organizzazione mantiene responsabilita sulle conseguenze dell'uso
  del riconoscimento facciale invece di delegarle interamente al
  modello.
\item
  La revisione introduce costi operativi e richiede informazioni
  sufficienti a ricostruire l'esito automatico.
\item
  Le scelte successive dovranno definire il processo concreto senza
  trasformare questa Decision in una procedura dettagliata.
\end{itemize}

\subsection{D1 freeze note}\label{d1-freeze-note}

This document ends Phase D1 authoring. No Decision above has yet been
classified as fitting or violating the frozen construction candidate.
Phase D2 must analyze this exact corpus before any semantic rewrite. If
D2 identifies a genuine counterexample, the original D1 text should
remain available as the evidence that triggered the model correction.

\end{document}
